\section{Introduction}
\label{sec:introduction}

Ideas compete. This observation is as old as intellectual history itself, but
its formal treatment has remained surprisingly underdeveloped. While
evolutionary biology has produced a rich mathematical apparatus for modeling
competition between species and strategies
\citep{hofbauer1998evolutionary,nowak2006evolutionary,weibull1995evolutionary},
and while cultural evolution theory has adapted some of these tools to the
study of beliefs and norms
\citep{cavalli1981cultural,boyd2005origin,mesoudi2011cultural}, the
specific dynamics of \emph{ideological} competition --- the struggle between
large-scale systems of political, religious, and philosophical thought ---
has largely been left to historians and political theorists working without
formal models.

This is not for lack of metaphor. Since \citet{dawkins1976selfish} introduced
the concept of the meme, the analogy between biological and cultural
evolution has been widely explored. But the analogy breaks in places that
matter. Three structural differences between species and ideas demand
mathematical attention:

\begin{enumerate}[label=(\roman*)]
  \item \textbf{Dormancy.} Ideas do not die when their carriers die or
  abandon them. Christianity survived the fall of Rome. Marxism survived the
  collapse of the Soviet Union. An idea can persist in latency --- in
  libraries, in diaspora communities, in private sympathy --- and reawaken
  when conditions change. In the standard replicator equation, the boundary
  $x_i = 0$ is absorbing: a strategy driven to zero stays at zero forever.
  For ideas, this is empirically wrong.

  \item \textbf{Hybridization.} A rabbit cannot also be a fox, but a person
  can be a Christian Marxist, or a nationalist who believes in universal
  rights. When ideas cohabit in a population, they can merge to form new
  species: social democracy emerged from the synthesis of Marxist analysis
  and liberal institutions \citep{berman2006primacy}; liberation theology
  from Christianity and anti-colonial politics. Standard population dynamics
  does not accommodate speciation during competition.

  \item \textbf{Endogenous landscape modification.} The fitness of a
  biological species depends on a niche that changes slowly relative to
  evolutionary timescales. The fitness of an idea depends on a social
  environment --- institutions, economic conditions, media structures ---
  that the idea itself modifies when it achieves dominance. Liberalism
  creates markets; markets create the inequality that fuels anti-liberal
  reaction. The fitness landscape is not exogenous: it is partly constructed
  by the competitors themselves.
\end{enumerate}

These are not minor perturbations of ecological dynamics. They are
qualitatively different phenomena that produce qualitatively different
behavior: oscillations driven by self-undermining dominance, resurrection
from apparent extinction, and the creation of genuinely new competitors
during the struggle.

The present paper develops a formal framework that addresses all three
features. We extend the replicator equation
\citep{taylor1978evolutionary} in three corresponding directions.
First, we introduce a \emph{latent sympathy} variable $y_i(t)$ alongside the
manifest adoption share $x_i(t)$, producing a replicator--mutator system
\citep{komarova2004replicator,page2002unifying} in which latent sympathy
prevents absorbing extinction and enables re-emergence. Suppression of an
ideology increases its latent sympathy (the \emph{martyrdom effect}),
formalizing the historical observation that persecuted ideas gain moral
authority. Second, we couple the replicator dynamics to an environment vector
$e(t)$ whose evolution depends on the current ideological distribution,
producing a co-evolutionary system with non-stationary fitness landscapes.
Third, we overlay a stochastic hybridization process that creates new
ideological types when latent sympathy overlap exceeds a threshold.

The resulting system is a piecewise-deterministic Markov process (PDMP)
\citep{davis1995piecewise}: continuous replicator--mutator dynamics
punctuated by discrete stochastic events (hybridization, mutation, charismatic
transmission). We implement this system in a simulation engine and apply it
to a historical scenario: four ideological currents --- liberalism, Marxism,
religious traditionalism, and nationalism --- competing in Europe from 1800 to
2025, subject to exogenous shocks calibrated to major historical crises.

The simulation is not a prediction engine. It is a tool for testing whether a
set of hypothesized mechanisms, when combined and allowed to interact, can
reproduce the qualitative shape of ideological history. We find that the
framework reproduces several robust features: the persistence of
traditionalism despite sustained modernization pressure, the crisis-driven
surges of radical movements \citep{funke2016going}, the dormancy and
reawakening of suppressed ideas, and the slow ideological drift produced by
endogenous environment modification.

We are explicit about the framework's deepest limitation. The replicator
equation, even in its extended form, treats individuals as passive carriers
of ideas whose adoption probabilities are determined by fitness and
environment. It has no room for genuine deliberation --- for a person who
sits quietly and changes their mind because they have thought carefully.
This is not an oversight; it is a structural choice that makes the population-level
dynamics tractable. But it means the model describes the ecology of ideas
\emph{as seen from above}. The view from inside --- from the deliberating
individual --- requires a different mathematics, perhaps one grounded in
decision theory or phenomenology rather than population dynamics. We regard
this boundary as the most important open problem in the field.

The paper is organized as follows. Section~\ref{sec:related} surveys related
work in cultural evolution, memetics, and computational social science.
Section~\ref{sec:framework} presents the mathematical framework.
Section~\ref{sec:software} describes the simulation software.
Section~\ref{sec:experiments} details the experimental design --- the choice
of ideological currents, parameter calibration, and historical shock
schedule. Section~\ref{sec:results} analyzes the simulation results.
Section~\ref{sec:discussion} discusses the contributions, limitations, and
implications of the framework.
